\documentclass[10pt]{article}

% Specify the margins and text size.
\setlength{\textwidth}{6.4in}
\setlength{\textheight}{9.5in}
\setlength{\oddsidemargin}{0pt}
\setlength{\evensidemargin}{0pt}
\setlength{\topmargin}{0pt}
\setlength{\hoffset}{.05in}
\setlength{\voffset}{-.8in}

\setlength{\parskip}{5pt}
\setlength{\parindent}{0pt}

% Load some fonts and symbol packages
\usepackage{latexsym}
\usepackage{pifont}       % contains 'star' symbol for counterinsurgency handbook title
\usepackage{yfonts} 
\usepackage{amsmath}
\usepackage{amsfonts}

\usepackage{graphicx}     % actually, this is loaded by pstricks
\usepackage[T1]{fontenc}
\usepackage{ifthen}
\usepackage{pstricks,pst-func,pst-text,pst-node,multido,pst-plot,calc,pst-3dplot}
%\usepackage[all]{xy}
%\usepackage{animate}

% The hyperref package inserts links into the pdf.
\definecolor{MyLinkColor}{rgb}{.1,.2,1}
\definecolor{MyCiteColor}{rgb}{.1,1,.2}
\definecolor{MyURLColor}{rgb}{.4,.4,.4}
\usepackage[backref=true,pagebackref=false,hyperindex,colorlinks=true,
  linkcolor=MyLinkColor,urlcolor=MyURLColor]{hyperref}


% The tweaklist package is something I found on the web.  It provides a simple interface
% for making changes to spacing used in the itemize and enumerate environments.  Comment
% this out if you don't care to use tweaklist.
\usepackage{tweaklist}
\renewcommand{\itemhook}{\setlength{\parskip}{2pt}\setlength{\parsep}%
{1pt}\setlength{\topsep}{0pt}\setlength{\itemsep}{0pt}}

\newcommand{\N}{\ensuremath\mathbb{N}}
\newcommand{\Q}{\ensuremath\mathbb{Q}}
\newcommand{\R}{\ensuremath\mathbb{R}}
\newcommand{\C}{\ensuremath\mathbb{C}}
\newcommand{\z}{\hphantom{0}}

\begin{document}
\pagestyle{empty}

\href{http://www.su.edu}{\includegraphics[height=1.75cm]{sulogo.eps}}
\vspace{-1.69cm}

{\small
\begin{tabular}{cl}
\hspace{5in} & Math 207:  Statistics\\
                & Discussion Exercises 1\\
\end{tabular}
}

\setlength{\baselineskip}{1.05\baselineskip}
\bigskip\bigskip


1. In the U.S. in 2000, there were 2.4 million deaths from all causes, compared to 1.9 million in 
197 --- a 25\% increase.  True or false, and explain:  the data show that the public's health got 
worse over the period 1970--2000.
\vspace{.5in}

2. Polio is an infectious disease;  for example, it seemed to spread when children went swimming together.  
The NFIP study was not done blind:  could that bias the results?  Discuss briefly.
\vspace{.5in}

3. Some studies find an association between liver cancer and smoking.  However,
alcohol consumption is a confounding variable.  This means ---

\hspace{10pt} i) Alcohol causes liver cancer.

\hspace{10pt} ii) Drinking is associated with smoking, and alcohol causes liver cancer.

Choose one and explain briefly.
\vspace{.5in}

4. The Federal Bureau of Investigation reports state-level and national data on crimes

\begin{itemize}
\item An investigator compares the incidence of crime in Minnesota and in Michigan.  
In 2001, there were 3,584 crimes in Minnesota, compared to 4,082 in Michigan.  
He concludes that Minnesotans are more law-abiding.  After all, Michigan includes the big bad
city of Detroit.  What do you say?\vspace{.5in}
%
\item An investigator compares the incidence of crime in the U.S. in 1991 and 2001.  In 1991, there
were 28,000 crimes, compared to 22,000 in 2001.  She concludes that
the U.S. became more law-abiding over that time period.  What do you say?
\end{itemize}
\vspace{.5in}

5. The National Highway and Traffic Safety Administration analyzed thefts of new cars in 2002, as well 
as sales figures for that year.

\begin{itemize}
\item There were 99 Corvettes stolen, and 26 Infiniti Q45 sedans.  Should you conclude
that American thieves prefer American cars?  or is something missing from the equation?
\vspace{.45in}
%
\item There were 50 BMW 7-series cars stolen, compared to 146 in the 3-series.  Should you conclude 
that thieves prefer smaller cars, which are more
economical to run and easier to park?  Or is something missing from the equation?
\vspace{.45in}
%
\item There were 429 Liberty Jeeps stolen, compared to 207,991 sold, for a rate of 2 per 100,000.  
True or false and explain:  the rate is low because the denominator is large.
\end{itemize}

\vfill
\eject

6. Ads for ADT Security system claim:\vspace{-5pt}
\begin{quotation}
When you go on vacation, burglars go to work $\dots$ According to  FBI statistics, 
over 25\% of home burglaries occur between Memorial Day and Labor Day.\vspace{-5pt}
\end{quotation}
Do the statistics prove that burglars go to work when other people go on vacation?  
Answer yes or no, and explain briefly.
\vspace{.5in}

7. California is evaluating a new program to rehabilitate prisoners before their release;  the object is to reduce
the recidivism rate --- the percentage who will be back in prison within
two years of release.  The program involves several months of 
\lq\lq boot camp\rq\rq\ --- military-style basic training with very strict discipline.
Admission to the program is voluntary.  According to a prison spokesman, \lq\lq Those who complete
boot camp are less likely to return to prison than other inmates.\rq\rq\

\begin{itemize}
\item What is the treatment group in the prison spokesman's comparison?
What is the control group?
\vspace{.5in}
%
\item Is the prison spokesman's comparison based on an observational study  or
a randomized controlled experiment?
\vspace{.5in}
%
\item True or false:  the data show that boot camp worked.
\end{itemize}


\vfill
\eject
\end{document}
